\subsection{Complete first ordinary-unit and S-unit norm images}

\begin{statement}{E.101}{Proved}{The remaining first-layer scalar and all four unit norm images}
Keep the literal fields and generators of E.86 and E.97:
\[
 \begin{gathered}
 \zeta^2+\zeta+1=0,\quad \xi^3=\zeta,\quad
 \beta_1^3=17\zeta,\quad\beta_2^3=\beta_1,\\
 L_i=\mathbb Q(\zeta,\beta_i),\quad B_i=L_i(\xi),
 \quad c_1=\beta_1+1+3\zeta.
 \end{gathered}
 \tag{E101.1}
\]
Let \(S\) be the primes above \(3,17\), and let a hat
denote completion for multiplication by \(3^m\).
The two certified free ordinary-unit generators of \(L_1\)
used in E.64 and E.97 are, with \(b_0=\beta_1\),
\[
 \begin{aligned}
 e_1&=-18+3b_0+3b_0^2-\frac{18b_0^3+5b_0^4-2b_0^5}{17},\\
 e_2&=-18-5b_0+2b_0^2-\frac{b_0^3+8b_0^4+b_0^5}{17},
 \qquad \epsilon=e_1e_2.
 \end{aligned}
 \tag{E101.2}
\]
For the exact ordinary divisor defect of E97.2, with
\(\sigma(\beta_2)=\zeta\beta_2\), one has
\[
       \omega_L(\epsilon)=(0,2),\qquad
       \omega_L(\zeta)=(0,1).
 \tag{E101.3}
\]
These coordinates are the specified
\((a_2\bmod3,b_2)\) coordinates of
\(\operatorname{Cl}(L_2)\{3\}/\langle\mathfrak q_{3,2}\rangle\).

Here are complete norm-image tests, including both directions.
Write an ordinary completed \(L_1\)-unit as
\(u=\zeta^c e_1^a e_2^b\), with
\(c\in\mathbb Z/3\), \(a,b\in\mathbb Z_3\).
For an S-unit append the factors \(c_1^e\beta_1^f\),
with \(e,f\in\mathbb Z_3\). Then
\[
 \begin{aligned}
 u\in N_{L_2/L_1}\widehat{\mathcal O_{L_2}^{\times}}
       &\Longleftrightarrow a\equiv b\equiv c\pmod3,\\
 uc_1^e\beta_1^f\in
       N_{L_2/L_1}\widehat{\mathcal O_{L_2,S}^{\times}}
       &\Longleftrightarrow
         a\equiv b\equiv c\pmod3\quad\hbox{and}\quad e\equiv0\pmod3.
 \end{aligned}
 \tag{E101.4}
\]
Thus the ordinary image has free basis
\(e_1^3,\zeta e_1e_2\); the S-unit image has free basis
\(e_1^3,\zeta e_1e_2,c_1^3,\beta_1\).
The actual quotient coordinate maps are respectively
\[
 \begin{aligned}
 \zeta^c e_1^a e_2^b&\longmapsto(a-b,c-b)\pmod3,\\
 \zeta^c e_1^a e_2^b c_1^e\beta_1^f
       &\longmapsto(a-b,c-b,e)\pmod3.
 \end{aligned}
 \tag{E101.5}
\]
Their inverses on quotient classes send
\((r,s)\) to \([e_1^r\zeta^s]\), and
\((r,s,t)\) to \([e_1^r\zeta^s c_1^t]\).
In particular these quotients are \(C_3^2\) and \(C_3^3\).

There is also a complete calculation over \(B_1\).
Use the exact E.79 basis
\[
 \mathcal B=(u_1,\ldots,u_{15})
   =(\xi,f_1,\ldots,f_8,s_1,\ldots,s_6),\qquad
 \widehat{\mathcal O_{B_1,S}^{\times}}
   =\mathbb Z_3^{15}/9\mathbb Z_3\mathbf e_1.
 \tag{E101.6}
\]
The \(f_i\) are its ordered certified fundamental units,
and \(s_i\) its ordered principal S-divisor generators.
The certificate retains their literal eighteen-term
coefficient vectors in the primitive element
\(X=\sqrt[3]{17}-\xi\); its polynomial is
\(X^{18}-105X^{15}+4137X^{12}-81182X^9
 +1128687X^6-6822222X^3+28934443\).
Thus (E101.6) fixes actual field elements and the
coordinate map \(x\mapsto\prod_i u_i^{x_i}\).

The norm image from \(B_2\), in these coordinates, is
exactly the six-congruence lattice
\[
 \begin{aligned}
 x_1-x_2-x_4-x_{10}-x_{12}&\equiv0, &
 x_5-2x_6-x_{10}-2x_{11}&\equiv0,\\
 x_8-x_{10}-x_{11}&\equiv0, &
 x_9-x_{10}-x_{12}-x_{15}&\equiv0,\\
 x_{13}-x_{15}&\equiv0, &x_{14}-x_{15}&\equiv0
 \qquad(\bmod3).
 \end{aligned}
 \tag{E101.7}
\]
Restricting to ordinary units, by setting
\(x_{10}=\cdots=x_{15}=0\), gives the complete test
\[
 x_1-x_2-x_4\equiv0,\qquad x_5-2x_6\equiv0,
 \qquad x_8\equiv x_9\equiv0\pmod3.
 \tag{E101.8}
\]
The two actual norm quotients are therefore \(C_3^6\)
and \(C_3^4\), respectively. The quotient map is the
ordered list of the six expressions in (E101.7), or
the four in (E101.8); a section on quotient classes
uses the coordinate vectors
\((\mathbf e_1,\mathbf e_5,\mathbf e_8,
  \mathbf e_9,\mathbf e_{13},\mathbf e_{14})\),
or its first four vectors. In these sentences a
section means a map into the quotient of units by
actual norms, so reducing an exponent modulo three
is justified by the included-base cube norm.

Here are explicit inverse coordinates on the S-unit
image itself. Define \(y_i=x_i\) for
\(i=2,3,4,6,7,10,11,12,15\), and set
\[
 \begin{aligned}
 y_1&=(x_1-x_2-x_4-x_{10}-x_{12})/3,\\
 y_5&=(x_5-2x_6-x_{10}-2x_{11})/3,\\
 y_8&=(x_8-x_{10}-x_{11})/3,\\
 y_9&=(x_9-x_{10}-x_{12}-x_{15})/3,\\
 y_{13}&=(x_{13}-x_{15})/3,
 \qquad y_{14}=(x_{14}-x_{15})/3.
 \end{aligned}
 \tag{E101.9}
\]
Recover \(x\) by these six equations and the unchanged
nine coordinates. The sole relation \(9\mathbf e_1\)
becomes \(3\mathbf e_1\) in the \(y\) coordinates.
Consequently the S-unit norm image is explicitly
\(C_3\oplus\mathbb Z_3^{14}\); the same formulas with
the last six coordinates zero identify the ordinary
image with \(C_3\oplus\mathbb Z_3^8\).

Every nonempty norm fibre is a translate of its
actual norm-one subgroup. For either ordinary units
or S-units the compact group types of these subgroups are
\[
 \begin{array}{c|c}
 L_2/L_1& C_3\oplus\mathbb Z_3^6\\
 B_2/B_1& C_3\oplus\mathbb Z_3^{18}.
 \end{array}
 \tag{E101.10}
\]
The proof gives the actual ideal equation defining
an ordinary \(L_2\)-unit \(\eta\) of norm
\(\zeta e_1e_2\). Its residue at seventeen has
\(\bar\eta^3=\zeta^2\) and \(\bar\eta^{96}=\zeta\),
so the ordinary upper unit group has a noncube residue.
The comparison to the positional section norm system
is the actual field norm of E.46 and E.88.
The manuscript anchor is LM \(\S51.16\),
folios733v--736v, scans172--178.
\end{statement}

\begin{proof}
\emph{A literal reversed norm and the actual eigenline.}
The lower polynomial is \(b^6+17b^3+289=0\).
Let \(U_0,U_1,U_2\in L_1\) be given by the following
matrix, whose rows multiply \((1,b,\ldots,b^5)^t\):
\[
 \begin{gathered}
 (U_0,U_1,U_2)^t=\frac1{16578519}\mathcal U
       (1,b,b^2,b^3,b^4,b^5)^t,\\
 \mathcal U^t=\begin{pmatrix}
 -37103278413&161479643843&-551544947582\\
 25225148828&-70558353324&-108462373370\\
 -5743099426&-3192102308&26395251042\\
 -1117338456&5646166658&-53737480949\\
 415980086&-6279927666&14839765429\\
 -548768638&1260332509&-2202740637
 \end{pmatrix}.
 \end{gathered}
 \tag{E101.11}
\]
Exact polynomial reduction gives
\[
 U_0^3+\epsilon U_1^3+\epsilon^2U_2^3
          -3\epsilon U_0U_1U_2=b.
 \tag{E101.12}
\]
Thus \(U_0+U_1\delta+U_2\delta^2\), for
\(\delta^3=\epsilon\), is a literal reversed norm
preimage of \(b\). Both GP and the independent rational
Python implementation verify (E101.12) from its coefficients.
The discovery computation used generators of an
uncertified BNF in this auxiliary field. Its class-group
output is not an input to (E101.11)--(E101.12) or to
any subsequent assertion.

Define the following matrices over \(L_1\):
\[
 \begin{gathered}
 A=\begin{pmatrix}0&0&\epsilon\\1&0&0\\0&1&0\end{pmatrix},
 \quad S=\operatorname{diag}(1,\zeta,\zeta^2),\\
 B=(U_0I+U_1A+U_2A^2)S,\\
 A^3=\epsilon I,\quad B^3=bI,\quad BA=\zeta AB.
 \end{gathered}
 \tag{E101.13}
\]
The certificate verifies each identity. In \(L_2\), the
\(\beta_2\)-eigenline of \(B\) is one-dimensional and
has a unique vector \(v\) with \(v_3=1\).
The same eigenline calculation as in E97.9 gives
\[
       Av=w\sigma(v),\qquad w=v_2,\qquad Nw=\epsilon.
 \tag{E101.14}
\]
Here the third coordinate gives the displayed formula
for \(w\). The certificate retains every coefficient
of \(v\) and \(w\) in the absolute degree-eighteen
field and independently verifies the relative and
absolute norm identities.

There is a short exact divisor construction which
avoids factoring the large numerator of \(w\).
Let \(I\) be the fractional ideal generated by the
three entries of \(v\), and put \(D=I^{-1}\).
Since \(\epsilon\) is an ordinary unit, both \(A\)
and \(A^{-1}\) have integral entries. Therefore the
ideal generated by \(Av\) is exactly \(I\).
Equation (E101.14) implies
\[
            I=(w)\sigma(I),\qquad
            (w)=\frac{\sigma(D)}D.
 \tag{E101.15}
\]
This is also checked by equality of the full ideal
HNF matrices. Thus \([D]\) is the actual class
defining \(\omega_L(\epsilon)\), with its sign fixed.

\emph{The complete class evaluation.}
An exact ideal reduction yields \(D=(m)J\), with
\(m\in L_2^\times\), and
\(N J=70782607=7^2\cdot1444543\).
The retained full ideal factorization has these three rows:
\[
 \begin{array}{c|c|c|c|c}
 p&f(P/p)&v_P(J)&f(P/(P\cap L_1))&
       \operatorname{Art}_{\Gamma_L}(P\cap L_1)\\ \hline
 7&1&1&1&(2,1,0)\\
 7&1&1&1&(2,1,0)\\
 1444543&1&1&1&(8,0,1).
 \end{array}
 \tag{E101.16}
\]
Both rational primes are certified prime. The two rows
over seven denote the two actual prime ideals stored
as HNF matrices. Contraction is tested by extension
of each lower prime's integral-basis generators and
its valuation at the upper prime; no root or prime
identification is inferred from the rational prime alone.
The Hilbert three-class field in E.86 is abelian over
\(L_1\), so the Artin values are the indicated lower
values multiplied by the relative residue degrees.
This is precisely the norm/Frobenius compatibility
\cite[V, Proposition3.3, Theorem3.6 and Example3.9]{E_MilneCFT}.

The sum of the last column is \((3,2,1)\) in
\(\Gamma_L=C_9\oplus C_3^2\).
Apply the actual E97.13 map and its retained inverse:
\[
 C=\begin{pmatrix}1&-3&-3\\0&1&0\\1/3&-2/3&-2/3\end{pmatrix},
 \qquad
 C^{-1}=\begin{pmatrix}-2&0&9\\0&1&0\\-1&-1&3\end{pmatrix}.
 \quad [D]_3=(3,2)\in C_9\oplus C_3.
 \tag{E101.17}
\]
The source is the subgroup where the third row of
\(C\) is integral. The target orders, including its
redundant coordinate, are \((9,3,1)\).
Quotienting by the strongly ambiguous generator
\((3,0)\) gives \((0,2)\), proving (E101.3).

E.97 proves that every ordinary lower unit is a field
norm and that the ordinary norm quotient has order nine,
with the exact Hilbert90 divisor map used here.
It also evaluates \(\omega_L(\zeta)=(0,1)\).
Thus \(\zeta\epsilon\) has zero defect.
The group of lower units modulo cubes has dimension three
over \(\mathbb F_3\); its norm kernel consequently has
dimension one. The nonzero vector \((c,a,b)=(1,1,1)\)
spans it. This proves the first equivalence in (E101.4),
and the quotient map and inverse in (E101.5) follow by
dividing by \((\zeta e_1e_2)^b\).

The S-unit norm image contains \(\beta_1=N\beta_2\)
and \(c_1^3=N(c_1)\), in addition to the ordinary
image. The subgroup they generate has index \(27\)
in the full lower S-unit group. E86.13 proves that
the actual S-unit norm quotient has order \(27\).
This inclusion and equality of indices prove the
second equivalence of (E101.4), its displayed free
basis, and the second quotient inverse.

\emph{An actual ordinary upper-unit preimage.}
Let \(D_\zeta,w_\zeta\) be the exact ideal and element
of E.97, and let \(D,w\) be those of (E101.14)--(E101.15).
The class of \(DD_\zeta/\mathfrak q_{3,2}\) is zero
in the three-primary coordinates. This ideal class
is fixed by \(\sigma\). It is also zero in the entire
class group: the ambiguous-class formula gives
\[
 \left|\operatorname{Cl}(L_2)^{\langle\sigma\rangle}\right|
    =\frac{h(L_1)\,e_3e_{17}}
       {3\,[E_1:E_1\cap NL_2^\times]}
    =\frac{9\cdot3\cdot3}{3}=27.
 \tag{E101.18}
\]
Every lower unit is a field norm, as established in E.97,
so the denominator's unit index is one. There are no
real ramified places. This applies the full ordinary
ambiguous-class formula, including its proof
\cite[Theorem1]{E_LemmermeyerAmbiguous}.
Thus no prime-to-three component is being discarded
when concluding principality.

Choose a generator \(d\) of
\(DD_\zeta/\mathfrak q_{3,2}\). Then the literal map
\[
        \eta=ww_\zeta\,\frac d{\sigma(d)}
        \in\mathcal O_{L_2}^{\times},\qquad
        N\eta=\zeta e_1e_2
 \tag{E101.19}
\]
follows by cancellation in the full divisor group.
This defines a preimage by an exact principal-ideal
equation whose solvability has just been proved.
Changing \(d\) by an ordinary unit changes \(\eta\)
by the corresponding norm-one Hilbert90 element.
For the ordinary image in (E101.4), write
\(a=b+3s\), \(c=b\bmod3\); its preimage is then
\(\eta^b e_1^s\). Appending \(c_1^{3t}\beta_1^f\)
gives the S-unit preimage \(\eta^b e_1^s c_1^t\beta_2^f\).
These formulas are maps on the completed groups as well.
The residue computations of E.97 give
\(\bar e_1=\bar e_2=\zeta^2\) at seventeen.
The vertical extension is totally ramified with
unchanged residue field \(\mathbb F_{289}\), so its
unit norm reduces to cubing. Hence
\(\bar\eta^3=\zeta^2\) and
\(\bar\eta^{96}=(\zeta^2)^{32}=\zeta\).

\emph{The full B-unit defect through the actual horizontal norm.}
Let \(R\) be the coefficient matrix of
\(N_h:B_1^\times\to L_1^\times\), in the basis
(E101.6) and target basis \((\zeta,e_1,e_2,c_1,\beta_1)\):
\[
 R=\left(\begin{array}{*{15}{r}}
 1&2&0&2&1&1&0&0&2&1&1&1&2&2&2\\
 0&0&0&0&-1&-1&0&0&-1&1&2&2&1&0&-1\\
 0&0&0&0&2&-1&0&1&1&1&1&0&0&0&1\\
 0&0&0&0&0&0&0&0&0&3&3&3&1&1&1\\
 0&0&0&0&0&0&0&0&0&0&0&0&1&1&1
 \end{array}\right).
 \tag{E101.20}
\]
The first row is interpreted modulo three.
The field morphism is given by
\(b\mapsto\sqrt[3]{17}\xi\); the certificate stores
its integral-basis matrix \(E\), together with
\((E^tE)^{-1}E^t\), and checks that this projection
recovers every actual horizontal norm.
For every column, multiplying its three literal
\(h\)-conjugates and reconstructing the target from
\(R\) agrees up to the prime-to-three sign.
This proves that (E101.20) is the actual norm map.

For reference the six local norm characters of E.79 are
\[
 M=\left(\begin{array}{*{15}{r}}
 2&1&0&1&0&0&0&2&1&1&1&0&0&0&2\\
 2&1&0&1&0&0&0&0&0&1&0&1&1&1&1\\
 2&1&0&1&0&0&0&1&2&1&2&2&0&0&1\\
 2&1&0&1&0&0&0&1&1&2&2&0&1&0&1\\
 2&1&0&1&0&0&0&1&1&2&2&0&0&1&1\\
 1&2&0&2&0&0&0&2&2&1&1&0&0&0&1
 \end{array}\right)\quad(\bmod3).
 \tag{E101.21}
\]
Its rank is five. For \(H=\ker M\subset\mathbb F_3^{15}\),
the exact matrix calculation gives
\[
        R(H)=\{(c,a,a,0,0):c,a\in\mathbb F_3\}.
 \tag{E101.22}
\]
Both inclusions and rank two are checked on a retained
ten-column basis of \(H\). Thus for a field-norm
S-unit \(u\), the horizontal norm can be divided by
included-base cubes and the known \(\beta_1\)-norm to
leave the ordinary target \(\zeta^c\epsilon^a\).
The commuting norm square of E97.15 and the exact
class map of E97.17 then give
\[
 \begin{aligned}
 \operatorname{Cl}_S(B_2)\{3\}&\longrightarrow
       \operatorname{Cl}_S(L_2)\{3\},&z&\longmapsto(0,z),\\
 \omega_L(N_hu)&=(0,c+2a),&
 \omega_B(u)&=c+2a.
 \end{aligned}
 \tag{E101.23}
\]
This follows directly on divisors:
\(\operatorname{div}_S x=(\sigma-1)D_B\) is sent by
\(N_h\) to \((\sigma-1)N_hD_B\); stripping the known
norm factors changes no defect class.
The injection \(z\mapsto(0,z)\) fixes the sign and
identifies the final scalar exactly.

Consequently the full B-S-unit norm image is the
kernel modulo three of the seven-row matrix
\[
             \Phi=\begin{pmatrix}M\\R_1+2R_2\end{pmatrix},
             \qquad \operatorname{rank}_{\mathbb F_3}\Phi=6.
 \tag{E101.24}
\]
The displayed \(R_i\) denote rows, indexed from one.
The first six rows impose field-norm membership by
the cyclic Hasse norm theorem; the last imposes zero
in the exact E86.14 S-unit norm defect. This proves
both directions before any lattice calculation.

The integral kernel has the HNF basis recovered
by reversing (E101.9). Its determinant is \(3^6\),
and \(\Phi\) kills each column modulo three.
The index therefore proves equality with the full
kernel. The retained Smith matrices have diagonal
\(3^6,1^9\), with both unimodular inverses checked.
Equations (E101.7) and (E101.9) follow.
The analogous ordinary submatrix has rank four;
its kernel HNF has diagonal \(3^4,1^5\).
Every S-unit norm preimage of an ordinary target
is already ordinary: each vertical S-prime is
unique above its predecessor with residue degree
one, so its norm valuation is exactly its own
valuation. This proves (E101.8) for ordinary norms.

\emph{Completion and fibre maps.}
All four unit groups are finitely generated.
Inclusion of a lower unit has norm its cube;
hence all norm cokernels here are finite and killed
by three. Completion for multiplication by \(3^m\)
preserves the calculated cokernels. The completed
norm images are closed images of compact
\(\mathbb Z_3\)-modules, so the congruence tests and
their inverse maps hold with exactly the topology
in (E101.4)--(E101.9).
Translation by any norm preimage gives an actual
homeomorphism from the norm-one subgroup onto
its fibre. In the L tower the unit ranks are eight
and two, and in the B tower they are twenty-six
and eight. The norm has full target rank because
its image contains cubes. The torsion kernel is
\(\mu_3\) in both cases: cubing on \(\mu_3\) or
\(\mu_9\), respectively. Thus (E101.10) follows
by the structure theorem for finitely generated
\(\mathbb Z_3\)-modules. The same S-valuation argument
shows that the ordinary and S-unit norm-one
subgroups are equal, as actual subgroups of the
upper field, not merely abstractly isomorphic.

The portable certificate independently checks the
literal polynomial norms, both eigenvector equations,
the class-coordinate inverses, every one of the
\(3^{10}\) field-norm residue vectors, and the full
Smith and projection inverses. It rebuilds only
the previously certified lower BNF/genus inputs.
The separate bounded search for a numerical
generator in (E101.19) is retained with its actual
no-hit output; its absence supplies no additional
mathematical assertion. No value of the eventual
exponent \(h\) or the E.91 transvection is deduced
from these first-layer norm calculations.
\end{proof}
